
%(BEGIN_QUESTION)
% Copyright 2014, Tony R. Kuphaldt, released under the Creative Commons Attribution License (v 1.0)
% This means you may do almost anything with this work of mine, so long as you give me proper credit

Forklar hvorfor vann (H$_{2}$O) ikke antenner, til tross for at det består av hydrogen som er svært brennbart og oksygen som er kjent for å akselerere forbrenning kraftig. Når disse to stoffene finnes i hvert vannmolekyl, hvorfor er ikke vann eksplosivt??

\vskip 20pt \vbox{\hrule \hbox{\strut \vrule{} {\bf Forslag til sokratisk diskusjon} \vrule} \hrule}

\begin{itemize}
\item{} Skriv en balansert kjemisk ligning som viser elektrolyse av vann til hydrogengass og oksygengass. Angi om $\Delta H$-verdien for denne ligningen er positiv eller negativ, og hva denne størrelsen betyr.
\item{} Knytt svaret ditt til tilfellet karbondioksid (CO$_{2}$), som består av karbon (C -- et brennbart grunnstoff) og oksygen (O -- en forbrenningsfremmer).
\end{itemize}

\underbar{file i01199}
%(END_QUESTION)





%(BEGIN_ANSWER)

Svaret på dette spørsmålet bygger på forståelse av hva forbrenning er, og hvordan energi frigjøres i kjemiske reaksjoner som forbrenning mellom hydrogen (H) og oksygen (O).
 
%(END_ANSWER)





%(BEGIN_NOTES)

Forbrenning -- som alle andre eksoterme kjemiske prosesser -- innebærer dannelse av bindinger mellom atomer som tidligere var adskilt.  Atomer som er bundet sammen har lavere energitilstand enn de samme atomene når de er adskilt, og derfor frigjøres energi når bindingene dannes og atomene går fra høyere til lavere energitilstand.

\vskip 10pt

Grunnen til at vann ikke forbrenner, er at {\it det allerede har brent}.  Vannmolekylet ble dannet da to hydrogenatomer bandt seg til ett oksygenatom og frigjorde energi.  Det er ingen energi igjen å frigjøre, fordi hydrogen- og oksygenatomene allerede har nådd sin lavere energitilstand ved dannelsen av vann.

\vskip 10pt

Det er bare når disse atomene finnes adskilt fra hverandre at det finnes potensial for forbrenning (energifrigjøring).

%INDEX% Kjemi, grunnleggende: molekylbindinger og energiutveksling

%(END_NOTES)

